\documentclass{encori}
\usepackage[utf8]{inputenc}
\usepackage[T1]{fontenc}
\usepackage{geometry}
\usepackage{graphicx}
\usepackage{float}
\usepackage{hyperref}
\usepackage[english,indonesian]{babel}

\title{Optimasi Pemrosesan Template LaTeX Menggunakan Pendekatan Berbasis AI}
\author{Academic Support Agent}
\date{\today}

\begin{document}
\maketitle

\begin{abstract}

\end{abstract}

\chapter{PENDAHULUAN}


\chapter{TINJAUAN PUSTAKA}


% Original content (commented out):
% \documentclass{encori}

%%%%%%%%%%%%%%%%%%%%%%%%%%%%%%%%%%%%%%%%%%%%%%%%%%%%%%%%%%%%%%%%%%%%%%%%
%% POR FAVOR, NÃO FAÇA MUDANÇAS NESSE PADRÃO QUE ACARRETEM  EM
%% ALTERAÇÃO NA FORMATAÇÃO FINAL DO TEXTO
%%%%%%%%%%%%%%%%%%%%%%%%%%%%%%%%%%%%%%%%%%%%%%%%%%%%%%%%%%%%%%%%%%%%%%%%

%%%%%%%%%%%%%%%%%%%%%%%%%%%%%%%%%%%%%%%%%%%%%%%%%%%%%%%%%%%%%%%%%%%%%%%%
% POR FAVOR, ESCOLHA CONFORME O CASO
%%%%%%%%%%%%%%%%%%%%%%%%%%%%%%%%%%%%%%%%%%%%%%%%%%%%%%%%%%%%%%%%%%%%%%%%
\usepackage[brazil]{babel} % texto em português
% \usepackage[english]{babel} % texto em inglês

%\usepackage[latin1]{inputenc} % acentuação em Português ISO-8859-1
\usepackage[utf8]{inputenc} % acentuação em Português UTF-8
%%%%%%%%%%%%%%%%%%%%%%%%%%%%%%%%%%%%%%%%%%%%%%%%%%%%%%%%%%%%%%%%%%%%%%%%


%%%%%%%%%%%%%%%%%%%%%%%%%%%%%%%%%%%%%%%%%%%%%%%%%%%%%%%%%%%%%%%%%%%%%%%%
%% POR FAVOR, NÃO ALTERAR
%%%%%%%%%%%%%%%%%%%%%%%%%%%%%%%%%%%%%%%%%%%%%%%%%%%%%%%%%%%%%%%%%%%%%%%%
\usepackage[T1]{fontenc}
\usepackage{float}
\usepackage{graphics}
\usepackage{graphicx}
\usepackage{epsfig}
\usepackage{indentfirst}
\usepackage{amsmath, amsfonts, amssymb, amsthm}
\usepackage{url}
\usepackage{csquotes}
\setlength{\parindent}{0em}
% Ambientes pré-definidos
\newtheorem{theorem}{Theorem}[section]
\newtheorem{lemma}{Lemma}[section]
\newtheorem{proposition}{Proposition}[section]
\newtheorem{definition}{Definition}[section]
\newtheorem{remark}{Remark}[section]
\newtheorem{corollary}{Corollary}[section]
\newtheorem{teorema}{Teorema}[section]
\newtheorem{lema}{Lema}[section]
\newtheorem{prop}{Proposi\c{c}\~ao}[section]
\newtheorem{defi}{Defini\c{c}\~ao}[section]
\newtheorem{obs}{Observa\c{c}\~ao}[section]
\newtheorem{cor}{Corol\'ario}[section]
\newcommand{\D}{\ensuremath{\mathbb{D}}}
\newcommand{\R}{\ensuremath{\mathbb{R}}} 
\newcommand{\Z}{\ensuremath{\mathbb{Z}}}
\newcommand{\N}{\ensuremath{\mathbb{N}}}
\newcommand{\NN}{\ensuremath{\mathcal{N}}}
\newcommand{\Q}{\ensuremath{\mathbb{Q}}}
\newcommand{\C}{\ensuremath{\mathbb{C}}}
\newcommand{\K}{\ensuremath{\mathbb{K}}}
\newcommand{\LL}{\ensuremath{\mathbb{L}}}
\newcommand{\Ll}{\ensuremath{\mathcal{L}}}
\newcommand{\F}{\ensuremath{\mathbb{F}}}
\newcommand{\UU}{\ensuremath{\mathbb{U}}}
\newcommand{\U}{\ensuremath{\mathcal{U}}}
\newcommand{\PP}{\ensuremath{\mathcal{P}}}
\newcommand{\A}{\ensuremath{\mathcal{A}}}
\newcommand{\Pp}{\ensuremath{\mathbb{P}}}
\newcommand{\G}{\ensuremath{\mathcal{G}}}
\newcommand{\OO}{\ensuremath{\mathcal{O}}}

% ref bibliográficas
\usepackage[backend=biber, style=numeric-comp, sorting=none]{biblatex}
\addbibresource{refs.bib}
\DeclareTextFontCommand{\emph}{\boldmath\bfseries}
\DefineBibliographyStrings{brazil}{phdthesis = {Tese de doutorado}}
\DefineBibliographyStrings{brazil}{mathesis = {Dissertação de mestrado}}
\DefineBibliographyStrings{english}{mathesis = {Master dissertation}}
%%%%%%%%%%%%%%%%%%%%%%%%%%%%%%%%%%%%%%%%%%%%%%%%%%%%%%%%%%%%%%%%%%%%%%%%


\begin{document}

%%%%%%%%%%%%%%%%%%%%%%%%%%%%%%%%%%%%%%%%%%%%%%%%%%%%%%%%%%%%%%%%%%%%%%%%
% TÍTULO E AUTORES
%%%%%%%%%%%%%%%%%%%%%%%%%%%%%%%%%%%%%%%%%%%%%%%%%%%%%%%%%%%%%%%%%%%%%%%%
\title{Título do Pôster/Apresentação}

% Se necessário, você deve acrescentar ou remover autores no comando abaixo. Se os autores pertencerem a mesma instituição é possível colocar apenas uma vez a instituição.

\author{
    {\large Autor 1}\thanks{autor1@email}\\ 
    {\small \textit{Instituição do(a) autor(a) principal}} \\
    {\large Autor 2}\thanks{autor2@email}\\
    {\small \textit{Instituição do(a) segundo(a) autor(a)}} 
}
\criartitulo
%%%%%%%%%%%%%%%%%%%%%%%%%%%%%%%%%%%%%%%%%%%%%%%%%%%%%%%%%%%%%%%%%%%%%%%%

%%%%%%%%%%%%%%%%%%%%%%%%%%%%%%%%%%%%%%%%%%%%%%%%%%%%%%%%%%%%%%%%%%%%%%%%
% RESUMO
%%%%%%%%%%%%%%%%%%%%%%%%%%%%%%%%%%%%%%%%%%%%%%%%%%%%%%%%%%%%%%%%%%%%%%%%

Este é o formato \LaTeX{} para a submissão de pôster e apresentação oral do \emph{EnCoRI}. \emph{Os trabalhos podem ser submetidos em português ou inglês, em forma de resumo de, no máximo, duas páginas, incluindo-se as referências bibliográficas.} Os resumos submetidos \emph{que não estiverem de acordo com o formato} apresentado por esse padrão serão \emph{rejeitados} pela organização do evento, sem análise do conteúdo científico.

Equações inseridas no resumo devem ser enumeradas à direita no texto, por exemplo,
\begin{equation}
y = x + z.
\label{eq1}
\end{equation}
Consulte o arquivo \verb!.tex! para acessar o código-fonte gerador da equação~\eqref{eq1}.

Sendo o resumo um texto curto, deve-se evitar o uso de tabelas e figuras. Se necessário, estas devem ser inseridas com o ambiente \verb!table!, conforme exemplo da Tabela~\ref{tabela01}.

\begin{table}[H]
\caption{{\small Exemplo de tabela.}}
\centering
\begin{tabular}{ccc}
\hline
Coluna 1  & Coluna 2 & Coluna 3\\ \hline
texto   & texto   & texto    \\
texto   & texto   & texto \\
\hline
\end{tabular}\label{tabela01}
\end{table}

Para inserir uma figura, deve-se usar o ambiente \verb!figure!, conforme exemplo da Figura~\ref{figura01}.

\begin{figure}[H]
\centering
\includegraphics[width=.3\textwidth]{exfig1}
\caption{{\small Exemplo de figura.}}
\label{figura01}
\end{figure}

As referências bibliográficas serão automaticamente geradas pelo modelo em ordem alfabética pelo sobrenome do primeiro autor. Este {\it template} fornece suporte para a inserção de referências bibliográficas com o Bib\LaTeX{}. Os dados de cada referência do resumo devem ser adicionados em \verb+refs.bib+ e a indicação da referência no texto deve ser inserida com o comando \verb+\cite+. Seguem alguns exemplos de referências: livro \cite{conway2013sphere}, artigos publicados em periódicos \cite{damir2021well}, dissertação de mestrado \cite{dantas2020funccoes}, tese de doutorado \cite{pinheiroorientador}, livro publicado dentro de uma série \cite{hefez2008codigos}, trabalho publicado em anais de eventos \cite{tonincodificaccao}, {\it website} e outros \cite{EnCoRI}. Por padrão deste modelo, as obras com mais de três autores, a referência irá indicar apenas o nome do primeiro autor, seguido da expressão et al.

\vspace{1.5em}

{\bf Palavras-chave:} De 2 a 5 palavras-chave separadas por vírgula.

\section*{Agradecimentos} % Você pode excluir a seção se preferir.
Caso necessário você pode fazer aqui seus agradecimentos, por exemplo, a agências de fomentos: CAPES, CNPq (314441/2021-2).


%%%%%%%%%%%%%%%%%%%%%%%%%%%%%%%%%%%%%%%%%%%%%%%%%%%%%%%%%%%%%%%%%%%%%%%%
% REFS BIBLIOGRÁFICAS
% POR FAVOR, NÃO ALTERAR
%%%%%%%%%%%%%%%%%%%%%%%%%%%%%%%%%%%%%%%%%%%%%%%%%%%%%%%%%%%%%%%%%%%%%%%%
\printbibliography
%%%%%%%%%%%%%%%%%%%%%%%%%%%%%%%%%%%%%%%%%%%%%%%%%%%%%%%%%%%%%%%%%%%%%%%%

\end{document}






\end{document}